\section{Approach}
\label{sec:approach}

Figure~\ref{fig:overview} shows the steps of \tool. The main idea is to
answer every question about a request from the specification first, and from
the API's earlier answers second. \tool reads the specification once (Steps 1
and~2). It then builds each request on demand (Steps 3 and~4). The engine
sends the request and stores the ids from 2XX responses in the resource pool,
so later requests can use real ids. Only Step~3 and Step~4 run per request,
and both run in microseconds (Section~\ref{sec:rq4}). The expensive part is
the HTTP round trip, which is not ours.

\begin{figure}[t]
  \centering
  %% Figure: overview of SpecSeed. Steps 1-4 are ours; the engine and the pool
%% belong to the team's engine. Grayscale-safe, \scriptsize, numbered steps.
\begin{tikzpicture}[
    x=1cm, y=1cm, font=\scriptsize,
    box/.style={draw, rounded corners=2pt, align=center, inner sep=2pt,
                minimum height=10mm, text width=17mm, fill=white},
    ext/.style={draw, dashed, rounded corners=2pt, align=center, inner sep=2pt,
                minimum height=10mm, text width=15mm, fill=black!6},
    flow/.style={-{Stealth[length=1.6mm]}, line width=0.5pt},
    back/.style={-{Stealth[length=1.6mm]}, line width=0.5pt, dash pattern=on 2pt off 1.5pt},
  ]
  \node[box, text width=12mm] (spec) at (0,0) {OpenAPI\\specification};
  \node[box] (s1) at (2.2,0) {\textbf{1} Parse and\\resolve};
  \node[box] (s2) at (4.6,0) {\textbf{2} Order\\operations};
  \node[box] (s3) at (7.0,0) {\textbf{3} Generate\\values};
  \node[box] (s4) at (7.0,-1.7) {\textbf{4} Reuse\\ids};
  \node[ext] (eng) at (9.6,0) {engine\\sends request};
  \node[ext, text width=10mm] (api) at (11.8,0) {REST\\API};
  \node[ext] (pool) at (9.6,-1.7) {resource\\pool};
  \draw[flow] (spec) -- (s1);
  \draw[flow] (s1) -- node[above] {operations} (s2);
  \draw[flow] (s2) -- node[above] {order} (s3);
  \draw[flow] (s3) -- node[above] {request} (eng);
  \draw[flow] (eng) -- (api);
  \draw[back] (api.south) |- node[pos=0.75, below] {2XX response} (pool.east);
  \draw[flow] (pool) -- node[below] {lookup} (s4);
  \draw[flow] (s4) -- node[right] {real ids} (s3);
  \node[font=\tiny, align=left, anchor=north west] at (1.2,-1.0)
    {\textbf{modes:} \mode{VALID} (nominal),\\\mode{BOUNDARY}, \mode{INVALID} (error)};
\end{tikzpicture}

  \caption{Overview of \tool. Steps 1 and 2 run once per specification;
    Steps 3 and 4 run for every request. The engine and the resource pool
    (dashed, shaded) belong to the team's engine: \tool only asks the pool
    for ids. The dashed arrow is the loop that turns 2XX responses into
    real ids for later requests.}
  \label{fig:overview}
\end{figure}

\subsection{Step 1: Parse and Resolve}
\label{sec:step1}
\tool turns the specification into a list of operations whose schemas
contain no references. For each path and each HTTP method under it, it
builds one operation with (i)~its parameters; (ii)~its request body;
(iii)~one schema per response status. Five rules decide the details.
\begin{enumerate}
  \item \emph{Any reference.} A \texttt{\$ref} is a pointer into the same
    document. \tool follows it by walking the document key by key, so the
    same code resolves schemas, parameters, request bodies and responses.
  \item \emph{Recursion.} \tool passes down the set of references it is
    currently inside. When it meets one of them again, it replaces that
    reference with an empty object schema and stops. The same reference used
    twice side by side is expanded twice, because it is not a cycle.
  \item \emph{Composition.} \texttt{allOf} becomes one object schema whose
    properties and required fields are the union of its parts.
    \texttt{oneOf} and \texttt{anyOf} become their first option.
  \item \emph{Parameters.} Parameters declared on the path apply to every
    method under it; a parameter of the operation with the same name and
    location replaces them. Path parameters are always required.
  \item \emph{Base path.} If the server URL has a path, such as
    \texttt{/api/v2}, \tool puts it in front of every operation path. The
    host and port come from the test setup, not from the specification.
\end{enumerate}
We do not claim this step as new: OpenAPI libraries resolve references too.
We wrote it ourselves to keep the component free of third-party packages
inside the competition container, and we tested it with 13 small
specifications, one rule each, before running it on real ones.

\subsection{Step 2: Order Operations}
\label{sec:step2}
\tool decides which operations create the resources that others need, and
sorts the operations so that creators come first. It names every resource
with a \emph{resource key}: a lower-case word without the plural and with
\texttt{id} added, so \texttt{owners}, \texttt{ownerId} and
\texttt{owner\_id} all become \texttt{ownerid}. A bare path parameter such as
\texttt{\{id\}} takes the name of the path segment before it.

An operation \emph{produces} a key when it is a \texttt{POST} to a collection
path, such as \texttt{POST /owners}. It \emph{consumes} a key for each path
parameter, and for each required body field that names another resource,
either a field ending in \texttt{Id} or an object with an \texttt{id}
property. When no operation produces a consumed key exactly, \tool takes the
shortest producer key that ends with it or shares its first six characters,
so the body field \texttt{type} (key \texttt{typeid}) links to
\texttt{POST /pettypes} (key \texttt{pettypeid}). The \emph{depth} of an operation is the length of the
longest chain of producers before it. The final order is: creators and
readers by increasing depth, and \texttt{DELETE} operations last, deepest
first, so a pet is deleted before its owner. The idea of producer--consumer
order comes from RESTler~\cite{restler} and the operation dependency graph of
RestTestGen~\cite{resttestgen}; we use a small name-based version of it.

\subsection{Step 3: Generate Values}
\label{sec:step3}
For each request, \tool fills every required parameter, about half of the
optional ones, and the request body. It has three modes. In \mode{VALID}
mode it tries to satisfy the schema and the API. For each value it takes the
first source that applies:
(i)~an \texttt{example} or \texttt{default} in the schema;
(ii)~an \texttt{enum} member;
(iii)~a value built for the \texttt{format}, such as an ISO date for
\texttt{date} or a well-formed UUID;
(iv)~a value built for the \texttt{pattern};
(v)~a value chosen from the name, such as an e-mail address for
\texttt{email}, ten digits for \texttt{telephone}, or a small integer for a
field ending in \texttt{Id};
(vi)~a random value of the right type inside the stated limits.
For objects it leaves out read-only fields, which the server sets.

\paragraph{Patterns.} To satisfy a \texttt{pattern}, \tool parses the regular
expression with Python's own regular-expression parser and walks the parse
tree: it emits literals, picks a character from each class, and repeats each
part a random number of times inside its bounds, preferring the minimum. It
then checks the result against the pattern and the length limits, and falls
back to simpler guesses (digits only, letters only) when the check fails.

\paragraph{Error modes.} In \mode{BOUNDARY} mode, numbers and lengths sit at
their stated limits. In \mode{INVALID} mode, \tool first builds a
\mode{VALID} request and then breaks exactly one thing: one parameter, one
body field, or the whole body (a wrong type, a value outside the limits, a
missing required field, or malformed JSON). The rest of the request stays
valid, so it passes the API's first checks and the error reaches deeper
code. We do not claim error modes as new; RestTestGen also separates nominal
and error testing~\cite{resttestgen-stvr}.

\subsection{Step 4: Reuse Ids}
\label{sec:step4}
Before Step~3 builds a value for a path or query parameter, or for a body
field that looks like an id, \tool asks the resource pool for a real one:
\texttt{pool.lookup(name, schema, op)}. If the pool has one, \tool uses it.
For a body field that is itself a resource, such as the \texttt{type} of a
pet, it asks for a whole object with \texttt{pool.lookup\_object} and sends
an existing pet type. If the pool has nothing, Step~3 makes up a value. The
pool lives in the engine; \tool only reads from it.
